\documentclass[10pt,a4paper]{amsart}
\usepackage{amsmath,amssymb,amsthm,graphicx}
\usepackage{verbatim}
\newenvironment{reference}{\comment}{\endcomment}
%\newenvironment{reference}{}{}
\newenvironment{slogan}{\comment}{\endcomment}
\newenvironment{history}{\comment}{\endcomment}
% For commutative diagrams we use Xy-pic
\usepackage[all]{xy}

% We use 2cell for 2-commutative diagrams.
\xyoption{2cell}
% Theorem environments.
%
\theoremstyle{plain}
\newtheorem{theorem}[subsection]{Theorem}
\newtheorem{proposition}[subsection]{Proposition}
\newtheorem{lemma}[subsection]{Lemma}

\theoremstyle{definition}
\newtheorem{definition}[subsection]{Definition}
\newtheorem{example}[subsection]{Example}
\newtheorem{exercise}[subsection]{Exercise}
\newtheorem{situation}[subsection]{Situation}

\theoremstyle{remark}
\newtheorem{remark}[subsection]{Remark}
\newtheorem{remarks}[subsection]{Remarks}

\numberwithin{equation}{subsection}

% Macros
%
\def\lim{\mathop{\mathrm{lim}}\nolimits}
\def\colim{\mathop{\mathrm{colim}}\nolimits}
\def\Spec{\mathop{\mathrm{Spec}}}
\def\Hom{\mathop{\mathrm{Hom}}\nolimits}
\def\Ext{\mathop{\mathrm{Ext}}\nolimits}
\def\SheafHom{\mathop{\mathcal{H}\!\mathit{om}}\nolimits}
\def\SheafExt{\mathop{\mathcal{E}\!\mathit{xt}}\nolimits}
\def\Sch{\mathit{Sch}}
\def\Mor{\mathop{\mathrm{Mor}}\nolimits}
\def\Ob{\mathop{\mathrm{Ob}}\nolimits}
\def\Sh{\mathop{\mathit{Sh}}\nolimits}
\def\NL{\mathop{N\!L}\nolimits}
\def\CH{\mathop{\mathrm{CH}}\nolimits}
\def\proetale{{pro\text{-}\acute{e}tale}}
\def\etale{{\acute{e}tale}}
\def\QCoh{\mathit{QCoh}}
\def\Ker{\mathop{\mathrm{Ker}}}
\def\Im{\mathop{\mathrm{Im}}}
\def\Coker{\mathop{\mathrm{Coker}}}
\def\Coim{\mathop{\mathrm{Coim}}}

% Boxtimes
%
\DeclareMathSymbol{\boxtimes}{\mathbin}{AMSa}{"02}

%
% Macros for moduli stacks/spaces
%
\def\QCohstack{\mathcal{QC}\!\mathit{oh}}
\def\Cohstack{\mathcal{C}\!\mathit{oh}}
\def\Spacesstack{\mathcal{S}\!\mathit{paces}}
\def\Quotfunctor{\mathrm{Quot}}
\def\Hilbfunctor{\mathrm{Hilb}}
\def\Curvesstack{\mathcal{C}\!\mathit{urves}}
\def\Polarizedstack{\mathcal{P}\!\mathit{olarized}}
\def\Complexesstack{\mathcal{C}\!\mathit{omplexes}}
% \Pic is the operator that assigns to X its picard group, usage \Pic(X)
% \Picardstack_{X/B} denotes the Picard stack of X over B
% \Picardfunctor_{X/B} denotes the Picard functor of X over B
\def\Pic{\mathop{\mathrm{Pic}}\nolimits}
\def\Picardstack{\mathcal{P}\!\mathit{ic}}
\def\Picardfunctor{\mathrm{Pic}}
\def\Deformationcategory{\mathcal{D}\!\mathit{ef}}
\usepackage{cleveref}
\crefname{originallemma}{original lemma}{original lemmas}
\Crefname{originallemma}{Original lemma}{Original lemmas}
\makeatletter
\newlabel{lemma-nonvanishing-generically}{{lemma-nonvanishing-generically}{source}{}{sourcepointer-1}{}}
\newlabel{lemma-nonvanishing-generically@cref}{{[originallemma][0][]lemma-nonvanishing-generically}{[1][0][]source}{}{}{}}
\makeatother
\crefname{originallemma}{original lemma}{original lemmas}
\Crefname{originallemma}{Original lemma}{Original lemmas}
\makeatletter
\newlabel{lemma-nonvanishing-generically-local}{{lemma-nonvanishing-generically-local}{source}{}{sourcepointer-2}{}}
\newlabel{lemma-nonvanishing-generically-local@cref}{{[originallemma][0][]lemma-nonvanishing-generically-local}{[1][0][]source}{}{}{}}
\makeatother
\crefname{originallemma}{original lemma}{original lemmas}
\Crefname{originallemma}{Original lemma}{Original lemmas}
\makeatletter
\newlabel{lemma-quotient-function}{{lemma-quotient-function}{source}{}{sourcepointer-3}{}}
\newlabel{lemma-quotient-function@cref}{{[originallemma][0][]lemma-quotient-function}{[1][0][]source}{}{}{}}
\makeatother
\crefname{originalreference}{original reference}{original references}
\Crefname{originalreference}{Original reference}{Original references}
\makeatletter
\newlabel{topology-definition-dimension-function}{{topology-definition-dimension-function}{source}{}{sourcepointer-4}{}}
\newlabel{topology-definition-dimension-function@cref}{{[originalreference][0][]topology-definition-dimension-function}{[1][0][]source}{}{}{}}
\makeatother
\title{The degree function attached to a dualizing complex is a dimension function}
\author{The Stacks Project Authors}
\date{stacks-project:dualizing.tex:local-snapshot-20261012; GNU FDL 1.2 or later}
\begin{document}
\maketitle
\noindent\textit{Editorial note.} All delivered mathematical text is the author's own, copied from the pinned Stacks Project chapter source sources/stacks-dualizing/pinned-original-source.tex; no statement or proof step was rewritten, reordered or removed.\par\medskip
\noindent\textit{Editorial note.} Preamble provenance: the Stacks Project keeps one single shared preamble for every chapter; this card takes it directly from the staged sources/stacks-dualizing/preamble.tex as three verbatim spans (lines 8-12, 14-18 and 154-217), so no derived file was written and nothing was added or redefined. The wrapper class and the base packages are supplied by the assembler.\par\medskip
\noindent\textit{Editorial note.} Bibliography provenance: neither the counted statement nor its proof contains a citation (this lane rescans the whole span before writing the card), so the card carries no bibliography and no entry was reconstructed.\par\medskip
\noindent\textit{Editorial note.} Selection provenance: this card was opened under the file-based control contract (CONTROL.md), with the mechanical pre-card gate run on its statement and proof interval before the card was written.\par\medskip
\noindent\textit{Allowed standard prerequisites (editorial).}\begin{itemize}
\item The reduction of the situation along a quotient by a prime, keeping a dualizing complex. Invoked inside the counted proof (source lines 3508-3512). Reduces an immediate specialization to the case of a domain with the smaller prime equal to zero.
\item The uniqueness of the degree in which the localization at a minimal prime is nonzero. Invoked inside the counted proof (source lines 3517-3518). Identifies the degree of the generic point of the localized complex.
\item The description of the degree function of a normalized dualizing complex over a local ring. Invoked inside the counted proof (source lines 3524-3527). Relates the degrees at the two primes through the dimension of the local ring at the larger prime.
\end{itemize}\medskip
\ifdefined\hypertarget\hypertarget{sourcepointer-1}{}\fi
\noindent\textit{Source reference.} Original lemma lemma-nonvanishing-generically: sources/stacks-dualizing/pinned-original-source.tex, source lines 3449--3449. Navigation pointer only: the counted text refers to the Stacks item lemma-nonvanishing-generically, which is an existing labeled result of this same chapter and is not part of this card. The pointer carries the target's own label as its number and its definition line as its span; its content is not reproduced here.\par\medskip
\ifdefined\hypertarget\hypertarget{sourcepointer-2}{}\fi
\noindent\textit{Source reference.} Original lemma lemma-nonvanishing-generically-local: sources/stacks-dualizing/pinned-original-source.tex, source lines 3403--3403. Navigation pointer only: the counted text refers to the Stacks item lemma-nonvanishing-generically-local, which is an existing labeled result of this same chapter and is not part of this card. The pointer carries the target's own label as its number and its definition line as its span; its content is not reproduced here.\par\medskip
\ifdefined\hypertarget\hypertarget{sourcepointer-3}{}\fi
\noindent\textit{Source reference.} Original lemma lemma-quotient-function: sources/stacks-dualizing/pinned-original-source.tex, source lines 3482--3482. Navigation pointer only: the counted text refers to the Stacks item lemma-quotient-function, which is an existing labeled result of this same chapter and is not part of this card. The pointer carries the target's own label as its number and its definition line as its span; its content is not reproduced here.\par\medskip
\ifdefined\hypertarget\hypertarget{sourcepointer-4}{}\fi
\noindent\textit{Source reference.} Original reference topology-definition-dimension-function: sources/stacks-dualizing/pinned-original-source.tex, source lines 3502--3502. Navigation pointer only: the counted text refers to the Stacks item topology-definition-dimension-function, which belongs to another chapter of the project and is not part of this card. The pointer carries the target's own label as its number because this chapter's staged snapshot does not contain the other chapter's tag table; its content is not reproduced here.\par\medskip
\begin{lemma}
\label{lemma-dimension-function}
Let $A$ be a Noetherian ring and let $\omega_A^\bullet$ be a dualizing
complex. The function $\delta = \delta_{\omega_A^\bullet}$
defined above is a dimension function
(Topology, Definition \ref{topology-definition-dimension-function}).
\end{lemma}

\begin{proof}
Let $\mathfrak p \subset \mathfrak q$ be an immediate specialization.
We have to show that $\delta(\mathfrak p) = \delta(\mathfrak q) + 1$.
We may replace $A$ by $A/\mathfrak p$, the complex $\omega_A^\bullet$ by
$\omega_{A/\mathfrak p}^\bullet = R\Hom(A/\mathfrak p, \omega_A^\bullet)$,
the prime $\mathfrak p$ by $(0)$, and the prime $\mathfrak q$
by $\mathfrak q/\mathfrak p$,
see Lemma \ref{lemma-quotient-function}. Thus we may assume that
$A$ is a domain, $\mathfrak p = (0)$, and $\mathfrak q$ is a prime
ideal of height $1$.

\medskip\noindent
Then $H^i(\omega_A^\bullet)_{(0)}$ is nonzero
for exactly one $i$, say $i_0$, by Lemma \ref{lemma-nonvanishing-generically}.
In fact $i_0 = -\delta((0))$ because
$(\omega_A^\bullet)_{(0)}[-\delta((0))]$
is a normalized dualizing complex over the field $A_{(0)}$.

\medskip\noindent
On the other hand $(\omega_A^\bullet)_\mathfrak q[-\delta(\mathfrak q)]$
is a normalized dualizing complex for $A_\mathfrak q$. By
Lemma \ref{lemma-nonvanishing-generically-local}
we see that
$$
H^e((\omega_A^\bullet)_\mathfrak q[-\delta(\mathfrak q)])_{(0)} =
H^{e - \delta(\mathfrak q)}(\omega_A^\bullet)_{(0)}
$$
is nonzero only for $e = -\dim(A_\mathfrak q) = -1$.
We conclude
$$
-\delta((0)) = -1 - \delta(\mathfrak q)
$$
as desired.
\end{proof}

\end{document}
